\section{Postulat Huntington dan Dasar Aksiomatik}

Struktur formal aljabar Boolean dibangun di atas seperangkat aksioma yang dikenal sebagai \logic{Postulat Huntington}, yang dirumuskan oleh matematikawan Amerika Edward Vermilye Huntington pada tahun 1904 \cite{huntington1904}. Postulat-postulat ini menetapkan syarat-syarat minimum yang harus dipenuhi oleh suatu sistem aljabar agar dapat dikategorikan sebagai aljabar Boolean. Perumusan aksiomatik ini penting karena memberikan fondasi rigorous yang memungkinkan kita membuktikan setiap teorema dan hukum aljabar Boolean secara deduktif, tanpa bergantung pada intuisi atau contoh spesifik semata.

Huntington menunjukkan bahwa hanya diperlukan enam postulat untuk mendefinisikan seluruh struktur aljabar Boolean secara lengkap. Dari keenam postulat ini, seluruh hukum dan teorema aljabar Boolean --- termasuk hukum asosiatif, hukum idempoten, hukum absorpsi, dan hukum De Morgan --- dapat diturunkan sebagai konsekuensi logis \cite{rosen2019}. Fakta bahwa hukum asosiatif \emph{bukan} merupakan postulat melainkan teorema turunan sering kali mengejutkan mahasiswa yang baru mempelajari aljabar Boolean, karena dalam aljabar konvensional, hukum asosiatif biasanya dianggap sebagai salah satu aksioma dasar.

\subsection{Enam Postulat Huntington}

Formulasi enam postulat berikut merupakan versi pedagogis yang umum dipakai di buku teks; Huntington (1904) menyajikan beberapa set aksioma yang ekuivalen dalam makalah aslinya \cite{huntington1904}. Berikut adalah keenam Postulat Huntington yang mendefinisikan aljabar Boolean $\langle B, +, \cdot, ', 0, 1 \rangle$:

\begin{enumerate}[label=\textbf{H\arabic*.}, itemsep=10pt]
  \item \textbf{Ketertutupan (\textit{Closure}):} Himpunan $B$ bersifat tertutup terhadap kedua operator biner. Artinya, untuk setiap $a, b \in B$, berlaku $a + b \in B$ dan $a \cdot b \in B$. Postulat ini menjamin bahwa hasil operasi pada elemen-elemen Boolean selalu menghasilkan elemen Boolean yang valid.

  \item \textbf{Elemen Identitas (\textit{Identity}):} Terdapat dua elemen identitas yang berbeda di dalam $B$:
  \begin{itemize}
      \item Elemen identitas untuk penjumlahan: $\exists\, 0 \in B$ sedemikian sehingga $a + 0 = a$ untuk setiap $a \in B$.
      \item Elemen identitas untuk perkalian: $\exists\, 1 \in B$ sedemikian sehingga $a \cdot 1 = a$ untuk setiap $a \in B$.
  \end{itemize}

  \item \textbf{Komutatif (\textit{Commutativity}):} Kedua operator biner bersifat komutatif. Untuk setiap $a, b \in B$:
  \[ a + b = b + a \quad \text{dan} \quad a \cdot b = b \cdot a \]

  \item \textbf{Distributif (\textit{Distributivity}):} Setiap operator bersifat distributif terhadap operator lainnya. Untuk setiap $a, b, c \in B$:
  \[ a \cdot (b + c) = (a \cdot b) + (a \cdot c) \quad \text{dan} \quad a + (b \cdot c) = (a + b) \cdot (a + c) \]

  \item \textbf{Komplemen (\textit{Complement}):} Untuk setiap elemen $a \in B$, terdapat elemen komplemen $a' \in B$ sedemikian sehingga:
  \[ a + a' = 1 \quad \text{dan} \quad a \cdot a' = 0 \]

  \item \textbf{Eksistensi Minimal (\textit{At Least Two Distinct Elements}):} Himpunan $B$ mengandung paling sedikit dua elemen yang berbeda, yaitu $0 \neq 1$.
\end{enumerate}

\begin{konsep}
Perhatikan bahwa Postulat Huntington memiliki sifat \textbf{dualitas}: untuk setiap postulat yang melibatkan operator $+$, terdapat postulat yang bersesuaian dengan operator $\cdot$, dan sebaliknya, dengan menukar $+$ dan $\cdot$ serta menukar $0$ dan $1$. Prinsip dualitas ini merupakan ciri khas aljabar Boolean yang membedakannya dari aljabar konvensional.
\end{konsep}

\subsection{Perbandingan dengan Aljabar Konvensional}

Aljabar Boolean memiliki beberapa perbedaan mendasar dengan aljabar konvensional (aljabar bilangan real), meskipun keduanya adalah struktur aljabar. Pemahaman terhadap perbedaan-perbedaan ini membantu mahasiswa menghindari kesalahan umum dalam manipulasi ekspresi Boolean. Tabel~\ref{tab:boolean-vs-konvensional} merangkum perbedaan utama antara kedua sistem aljabar tersebut.

\begin{table}[H]
\centering
\renewcommand{\arraystretch}{1.3}
\begin{tabular}{|p{4cm}|p{5cm}|p{5cm}|}
\hline
\textbf{Aspek} & \textbf{Aljabar Konvensional} & \textbf{Aljabar Boolean} \\ \hline
Domain variabel & Bilangan real ($\mathbb{R}$) & $\{0, 1\}$ \\ \hline
Operator dasar & $+, -, \times, \div$ & $+$ (OR), $\cdot$ (AND), $'$ (NOT) \\ \hline
Invers aditif & $a + (-a) = 0$ & Tidak ada invers aditif \\ \hline
Invers multiplikatif & $a \times a^{-1} = 1$ (untuk $a \neq 0$) & Tidak ada invers multiplikatif \\ \hline
Hukum distributif & Hanya $\times$ atas $+$ & Keduanya: $\cdot$ atas $+$ DAN $+$ atas $\cdot$ \\ \hline
Hukum komplemen & Tidak ada & $a + a' = 1$, $a \cdot a' = 0$ \\ \hline
Idempoten & Tidak berlaku: $a + a \neq a$ & Berlaku: $a + a = a$, $a \cdot a = a$ \\ \hline
\end{tabular}
\caption{Perbandingan Aljabar Konvensional dan Aljabar Boolean}
\label{tab:boolean-vs-konvensional}
\end{table}

Perbedaan paling signifikan terletak pada distributivitas ganda dan hukum komplemen. Dalam aljabar konvensional, distributivitas hanya berlaku untuk perkalian terhadap penjumlahan ($a \times (b + c) = ab + ac$), tetapi penjumlahan \emph{tidak} distributif terhadap perkalian ($a + (b \times c) \neq (a+b)(a+c)$ secara umum). Sebaliknya, dalam aljabar Boolean, kedua bentuk distributif berlaku secara simetris. Selain itu, konsep invers (pengurangan dan pembagian) tidak ada dalam aljabar Boolean --- satu-satunya ``operasi kebalikan'' yang tersedia adalah komplemen \cite{rosen2019}.

Tidak adanya operasi pengurangan dan pembagian dalam aljabar Boolean bukanlah kelemahan, melainkan konsekuensi alami dari sifat himpunan $\{0, 1\}$ yang hanya memiliki dua elemen. Dalam konteks rangkaian digital, keterbatasan ini justru menjadi keunggulan: setiap variabel hanya memiliki dua kemungkinan keadaan (hidup/mati, tinggi/rendah), sehingga manipulasi aljabar menjadi lebih sederhana dan dapat diimplementasikan secara langsung menggunakan transistor semikonduktor \cite{allaboutcircuits_boolean}. Kesederhanaan ini yang menjadikan aljabar Boolean menjadi fondasi ideal bagi seluruh teknik desain rangkaian digital modern.
